% Shared preamble of the TikZ schematics (Fig. 2, Fig. 3), lane p0-w6.
% Built by `nacre figs` (nacre.figs.schematic) with pdflatex through latexmk. The PDF omits its
% dates and trailer id so that a rebuild is byte-identical.
\ifdefined\pdftrailerid
  \pdfinfoomitdate=1 \pdftrailerid{} \pdfsuppressptexinfo=-1
\fi% XeLaTeX (symbol variants): dates and id are fixed by SOURCE_DATE_EPOCH at build time
\usetikzlibrary{arrows.meta,calc}
% Generated by nacre.figs.pictograms from figures/rdkit/structures.yaml and figures/style;
% do not edit. Text-free structure pictograms for the TikZ schematics, one bond length = 1.
% Strokes: applied.schematic_stroke (bonds, disc outlines), applied.schematic_stroke_bold (arrows).
\definecolor{elH}{rgb}{1.0000,1.0000,1.0000}
\definecolor{elC}{rgb}{0.5650,0.5650,0.5650}
\definecolor{elN}{rgb}{0.1880,0.3140,0.9730}
\definecolor{elO}{rgb}{1.0000,0.0510,0.0510}
\definecolor{elAl}{rgb}{0.8419,0.7566,0.6633}
\definecolor{elCa}{rgb}{0.2390,1.0000,0.0000}
\definecolor{okblack}{rgb}{0.0000,0.0000,0.0000}
\definecolor{okorange}{rgb}{0.9020,0.6235,0.0000}
\definecolor{okskyblue}{rgb}{0.3373,0.7059,0.9137}
\definecolor{okbluishgreen}{rgb}{0.0000,0.6196,0.4510}
\definecolor{okyellow}{rgb}{0.9412,0.8941,0.2588}
\definecolor{okblue}{rgb}{0.0000,0.4471,0.6980}
\definecolor{okvermillion}{rgb}{0.8353,0.3686,0.0000}
\definecolor{okreddishpurple}{rgb}{0.8000,0.4745,0.6549}
\definecolor{okgray}{rgb}{0.6000,0.6000,0.6000}
\tikzset{bond/.style={line width=0.5bp, line cap=round, draw=okblack},
  atom/.style={line width=0.5bp, draw=okblack},
  thin/.style={line width=0.5bp},
  bold/.style={line width=1.6bp}}
\tikzset{pics/catechol/.style={code={
  \draw[bond] (-0.5000,-1.2990) -- (0.5000,-1.2991);
  \draw[bond] (-0.3200,-1.0790) -- (0.3200,-1.0791);
  \draw[bond] (0.5000,-1.2991) -- (1.0000,-0.4330);
  \draw[bond] (1.0000,-0.4330) -- (0.5000,0.4331);
  \draw[bond] (0.7195,-0.3871) -- (0.3995,0.1672);
  \draw[bond] (0.5000,0.4331) -- (-0.5000,0.4329);
  \draw[bond] (-0.5000,0.4329) -- (-0.9999,-0.4330);
  \draw[bond] (-0.3995,0.1671) -- (-0.7194,-0.3871);
  \draw[bond] (-0.5000,0.4329) -- (-1.0000,1.2990);
  \draw[bond] (0.5000,0.4331) -- (1.0000,1.2991);
  \draw[bond] (-0.9999,-0.4330) -- (-0.5000,-1.2990);
  \coordinate (-a0) at (-0.5000,-1.2990);
  \coordinate (-a1) at (0.5000,-1.2991);
  \coordinate (-a2) at (1.0000,-0.4330);
  \coordinate (-a3) at (0.5000,0.4331);
  \coordinate (-a4) at (-0.5000,0.4329);
  \coordinate (-a5) at (-0.9999,-0.4330);
  \coordinate (-a6) at (-1.0000,1.2990);
  \filldraw[atom, fill=elO] (-1.0000,1.2990) circle (0.300);
  \coordinate (-a7) at (1.0000,1.2991);
  \filldraw[atom, fill=elO] (1.0000,1.2991) circle (0.300);
  \coordinate (-head) at (0.0000,1.2991);
}}}
\tikzset{pics/catecholatedianion/.style={code={
  \draw[bond] (-1.0000,1.2991) -- (-0.4999,0.4330);
  \draw[bond] (-0.4999,0.4330) -- (-1.0001,-0.4330);
  \draw[bond] (-0.3995,0.1671) -- (-0.7195,-0.3871);
  \draw[bond] (-1.0001,-0.4330) -- (-0.5000,-1.2990);
  \draw[bond] (-0.5000,-1.2990) -- (0.5000,-1.2990);
  \draw[bond] (-0.3200,-1.0790) -- (0.3200,-1.0790);
  \draw[bond] (0.5000,-1.2990) -- (0.9999,-0.4330);
  \draw[bond] (0.9999,-0.4330) -- (0.5001,0.4330);
  \draw[bond] (0.7194,-0.3871) -- (0.3995,0.1671);
  \draw[bond] (0.5001,0.4330) -- (1.0000,1.2990);
  \draw[bond] (0.5001,0.4330) -- (-0.4999,0.4330);
  \coordinate (-a0) at (-1.0000,1.2991);
  \filldraw[atom, fill=elO] (-1.0000,1.2991) circle (0.300);
  \coordinate (-a1) at (-0.4999,0.4330);
  \coordinate (-a2) at (-1.0001,-0.4330);
  \coordinate (-a3) at (-0.5000,-1.2990);
  \coordinate (-a4) at (0.5000,-1.2990);
  \coordinate (-a5) at (0.9999,-0.4330);
  \coordinate (-a6) at (0.5001,0.4330);
  \coordinate (-a7) at (1.0000,1.2990);
  \filldraw[atom, fill=elO] (1.0000,1.2990) circle (0.300);
  \coordinate (-head) at (-0.0000,1.2990);
}}}
\tikzset{pics/acetate/.style={code={
  \draw[bond] (-0.0000,-1.0000) -- (0.0000,-0.0000);
  \draw[bond] (-0.0550,-0.0953) -- (-0.9210,0.4048);
  \draw[bond] (0.0550,0.0952) -- (-0.8110,0.5953);
  \draw[bond] (0.0000,-0.0000) -- (0.8660,0.5000);
  \coordinate (-a0) at (-0.0000,-1.0000);
  \coordinate (-a1) at (0.0000,-0.0000);
  \coordinate (-a2) at (-0.8660,0.5000);
  \filldraw[atom, fill=elO] (-0.8660,0.5000) circle (0.300);
  \coordinate (-a3) at (0.8660,0.5000);
  \filldraw[atom, fill=elO] (0.8660,0.5000) circle (0.300);
  \coordinate (-head) at (0.0000,0.5000);
}}}
\tikzset{pics/oligomerunit/.style={code={
  \draw[bond] (-1.0000,1.9661) -- (-0.5001,1.1001);
  \draw[bond] (-0.5001,1.1001) -- (-1.0001,0.2340);
  \draw[bond] (-1.0001,0.2340) -- (-0.5000,-0.6320);
  \draw[bond] (-0.7196,0.1882) -- (-0.3995,-0.3661);
  \draw[bond] (-0.5000,-0.6320) -- (-0.8089,-1.5830);
  \draw[bond] (-0.8089,-1.5830) -- (0.0001,-2.1708);
  \draw[bond] (-0.5340,-1.5108) -- (-0.0162,-1.8870);
  \draw[bond] (0.0001,-2.1708) -- (0.8091,-1.5828);
  \draw[bond] (0.8091,-1.5828) -- (0.5000,-0.6319);
  \draw[bond] (0.5000,-0.6319) -- (0.9999,0.2341);
  \draw[bond] (0.3995,-0.3661) -- (0.7194,0.1882);
  \draw[bond] (0.9999,0.2341) -- (0.5000,1.1001);
  \draw[bond] (0.5000,1.1001) -- (1.0000,1.9662);
  \draw[bond] (0.5000,1.1001) -- (-0.5001,1.1001);
  \draw[bond] (0.3200,0.8801) -- (-0.3201,0.8801);
  \draw[bond] (0.5000,-0.6319) -- (-0.5000,-0.6320);
  \coordinate (-a0) at (-1.0000,1.9661);
  \filldraw[atom, fill=elO] (-1.0000,1.9661) circle (0.300);
  \coordinate (-a1) at (-0.5001,1.1001);
  \coordinate (-a2) at (-1.0001,0.2340);
  \coordinate (-a3) at (-0.5000,-0.6320);
  \coordinate (-a4) at (-0.8089,-1.5830);
  \coordinate (-a5) at (0.0001,-2.1708);
  \coordinate (-a6) at (0.8091,-1.5828);
  \filldraw[atom, fill=elN] (0.8091,-1.5828) circle (0.300);
  \coordinate (-a7) at (0.5000,-0.6319);
  \coordinate (-a8) at (0.9999,0.2341);
  \coordinate (-a9) at (0.5000,1.1001);
  \coordinate (-a10) at (1.0000,1.9662);
  \filldraw[atom, fill=elO] (1.0000,1.9662) circle (0.300);
  \coordinate (-head) at (0.0000,1.9661);
}}}
\tikzset{pics/chitosanunit/.style={code={
  \draw[bond] (-1.6547,-0.6785) -- (-0.7877,-0.1803);
  \draw[bond] (-0.7877,-0.1803) -- (0.0773,-0.6821);
  \draw[bond] (0.0773,-0.6821) -- (0.9443,-0.1837);
  \draw[bond] (0.9443,-0.1837) -- (0.9465,0.8163);
  \draw[bond] (0.9465,0.8163) -- (0.0815,1.3180);
  \draw[bond] (0.0815,1.3180) -- (-0.7857,0.8197);
  \draw[bond] (0.0815,1.3180) -- (0.0826,1.8680);
  \draw[bond] (0.9465,0.8163) -- (1.8135,1.3146);
  \draw[bond] (0.9443,-0.1837) -- (1.8095,-0.6855);
  \draw[bond] (0.0773,-0.6821) -- (0.0753,-1.6819);
  \draw[bond] (-1.6547,-0.6785) -- (-2.5197,-0.1769);
  \draw[bond] (-0.7857,0.8197) -- (-0.7877,-0.1803);
  \draw[bond] (0.0753,-1.6819) -- (-0.4015,-1.9560);
  \coordinate (-a0) at (-1.6547,-0.6785);
  \coordinate (-a1) at (-0.7877,-0.1803);
  \coordinate (-a2) at (0.0773,-0.6821);
  \coordinate (-a3) at (0.9443,-0.1837);
  \coordinate (-a4) at (0.9465,0.8163);
  \coordinate (-a5) at (0.0815,1.3180);
  \coordinate (-a6) at (-0.7857,0.8197);
  \filldraw[atom, fill=elO] (-0.7857,0.8197) circle (0.300);
  \coordinate (-a7) at (0.0835,2.3180);
  \coordinate (-a8) at (1.8135,1.3146);
  \filldraw[atom, fill=elN] (1.8135,1.3146) circle (0.300);
  \coordinate (-a9) at (1.8095,-0.6855);
  \filldraw[atom, fill=elO] (1.8095,-0.6855) circle (0.300);
  \coordinate (-a10) at (0.0753,-1.6819);
  \filldraw[atom, fill=elO] (0.0753,-1.6819) circle (0.300);
  \coordinate (-a11) at (-2.5197,-0.1769);
  \filldraw[atom, fill=elO] (-2.5197,-0.1769) circle (0.300);
  \coordinate (-a12) at (-0.7917,-2.1803);
}}}
\tikzset{pics/water/.style={code={
  \draw[bond] (0.0000,0.0000) -- (0.8660,-0.5001);
  \draw[bond] (0.0000,0.0000) -- (-0.8660,-0.5001);
  \coordinate (-a0) at (0.0000,0.0000);
  \filldraw[atom, fill=elO] (0.0000,0.0000) circle (0.300);
  \coordinate (-a1) at (0.8660,-0.5001);
  \filldraw[atom, fill=elH] (0.8660,-0.5001) circle (0.200);
  \coordinate (-a2) at (-0.8660,-0.5001);
  \filldraw[atom, fill=elH] (-0.8660,-0.5001) circle (0.200);
  \coordinate (-head) at (0.0000,0.0000);
}}}
\tikzset{pics/carbonate/.style={code={
  \draw[bond] (-0.0551,0.0952) -- (0.8105,0.5958);
  \draw[bond] (0.0551,-0.0952) -- (0.9207,0.4054);
  \draw[bond] (0.0000,-0.0000) -- (-0.8664,0.4994);
  \draw[bond] (0.0000,-0.0000) -- (0.0008,-1.0000);
  \coordinate (-a0) at (0.0000,-0.0000);
  \coordinate (-a1) at (0.8656,0.5006);
  \filldraw[atom, fill=elO] (0.8656,0.5006) circle (0.300);
  \coordinate (-a2) at (-0.8664,0.4994);
  \filldraw[atom, fill=elO] (-0.8664,0.4994) circle (0.300);
  \coordinate (-a3) at (0.0008,-1.0000);
  \filldraw[atom, fill=elO] (0.0008,-1.0000) circle (0.300);
}}}

\tikzset{
  proc/.style={bold, draw=okblack, -{Stealth[length=7bp,width=5.5bp]}},
  force/.style={bold, -{Stealth[length=6bp,width=4.5bp]}},
  statebox/.style={fill=okgray!14, draw=none, rounded corners=3bp},
  oxide/.style={fill=elAl!55, draw=okblack, thin},
  template/.style={fill=okorange!30, draw=none},
  nucleus/.style={fill=okskyblue!35, draw=none},
  ion/.style={atom, fill=elCa},
  placeholder/.style={draw=okblack, thin, dash pattern=on 2bp off 1.5bp, fill=white},
}
% Ca2+ ion disc and a hydrated Ca2+ (six waters, oxygen towards the ion).
\newcommand{\caion}[1]{\filldraw[ion] (#1) circle (0.16);}
\newcommand{\hydratedca}[2]{% #1 centre, #2 pictogram scale (cm per bond length)
  \foreach \a in {0,60,...,300} {\pic[scale=#2, rotate=\a+90] at ($(#1)+(\a:0.58)$) {water};}
  \caion{#1}}
% Oxide strip with surface hydroxyls, from x=#1 to x=#2, top at y=#3 (illustrative).
\newcommand{\oxidestrip}[3]{%
  \fill[oxide] (#1,#3-0.55) rectangle (#2,#3);
  \draw[thin] (#1,#3) -- (#2,#3);
  \pgfmathsetmacro{\oxa}{#1+0.2}\pgfmathsetmacro{\oxb}{#1+0.6}\pgfmathsetmacro{\oxe}{#2-0.1}%
  \foreach \x in {\oxa,\oxb,...,\oxe} {\filldraw[atom, fill=elO] (\x,#3) circle (0.075);}}
% A placeholder slot for a render supplied later (lane p0-w7): dashed frame with diagonals.
\newcommand{\placeholder}[4]{% lower left (#1,#2), upper right (#3,#4)
  \filldraw[placeholder] (#1,#2) rectangle (#3,#4);
  \draw[thin, okgray] (#1,#2) -- (#3,#4);
  \draw[thin, okgray] (#3,#2) -- (#1,#4);}
